\documentclass[12pt,a4paper]{article}
\usepackage[utf8]{inputenc}
\usepackage{amsmath,amssymb,amsthm}
\usepackage{algorithm}
\usepackage{algorithmic}
\usepackage{geometry}
\geometry{margin=1in}

\newtheorem{theorem}{Theorem}
\newtheorem{lemma}[theorem]{Lemma}
\newtheorem{proposition}[theorem]{Proposition}
\theoremstyle{definition}
\newtheorem{definition}{Definition}

\title{Project Euler Problem 21: Amicable Numbers}
\author{}
\date{}

\begin{document}
\maketitle

\section{Problem Statement}

Let $d(n)$ denote the sum of proper divisors of~$n$. If $d(a) = b$ and $d(b) = a$ with $a \neq b$, then $a$ and $b$ form an \emph{amicable pair}. Compute the sum of all amicable numbers below $10000$.

\section{Definitions}

\begin{definition}[Divisor-Sum Functions]
For $n \in \mathbb{Z}_{>0}$, define $\sigma_1(n) = \sum_{d \mid n} d$ and $s(n) = \sigma_1(n) - n$.
\end{definition}

\begin{definition}[Amicable Pair]
Distinct positive integers $a, b$ form an \emph{amicable pair} if $s(a) = b$ and $s(b) = a$.
\end{definition}

\section{Theorems and Proofs}

\begin{theorem}[Multiplicativity of $\sigma_1$]
If $\gcd(m, n) = 1$, then $\sigma_1(mn) = \sigma_1(m)\,\sigma_1(n)$. Consequently,
\[
\sigma_1\!\Bigl(\prod_{i=1}^{k} p_i^{a_i}\Bigr) = \prod_{i=1}^{k} \frac{p_i^{a_i+1} - 1}{p_i - 1}.
\]
\end{theorem}

\begin{proof}
The map $(d_1, d_2) \mapsto d_1 d_2$ is a bijection from $\{d_1 : d_1 \mid m\} \times \{d_2 : d_2 \mid n\}$ to $\{d : d \mid mn\}$, by unique factorization and coprimality. Hence
\[
\sigma_1(mn) = \sum_{d_1 \mid m}\;\sum_{d_2 \mid n} d_1 d_2 = \sigma_1(m)\,\sigma_1(n).
\]
For a prime power, $\sigma_1(p^a) = \sum_{j=0}^{a} p^j = (p^{a+1}-1)/(p-1)$ by the geometric series. The product formula follows by induction on the number of distinct prime factors.
\end{proof}

\begin{lemma}[Trial-Division Computation]
The value $s(n)$ can be computed in $O(\sqrt{n})$ arithmetic operations.
\end{lemma}

\begin{proof}
Divisors of $n$ pair as $(d, n/d)$ for $1 \leq d \leq \sqrt{n}$, so
\[
\sigma_1(n) = \sum_{\substack{d=1 \\ d \mid n}}^{\lfloor\sqrt{n}\rfloor} \!\bigl(d + n/d\bigr) - \epsilon(n), \quad
\epsilon(n) = \begin{cases} \sqrt{n} & \text{if } \lfloor\sqrt{n}\rfloor^2 = n, \\ 0 & \text{otherwise.}\end{cases}
\]
The loop runs $\lfloor\sqrt{n}\rfloor$ iterations at $O(1)$ each, giving $O(\sqrt{n})$ total. Then $s(n) = \sigma_1(n) - n$.
\end{proof}

\begin{theorem}[Enumeration]
The complete set of amicable pairs $(a,b)$ with $a < b < 10000$ is:
\[
\{(220, 284),\; (1184, 1210),\; (2620, 2924),\; (5020, 5564),\; (6232, 6368)\}.
\]
\end{theorem}

\begin{proof}
Exhaustive computation: for each $n \in [2, 9999]$, compute $m = s(n)$; if $m \neq n$ and $s(m) = n$, record $\{n, m\}$. This finite, deterministic enumeration yields exactly the five pairs listed.
\end{proof}

\section{Algorithm}

\begin{algorithm}[H]
\caption{Sum of Amicable Numbers below $N$}
\begin{algorithmic}[1]
\REQUIRE Upper bound $N$
\ENSURE Sum of all amicable numbers in $[2, N)$
\STATE $\mathrm{total} \gets 0$
\FOR{$n = 2$ \TO $N - 1$}
    \STATE $m \gets s(n)$
    \IF{$m \neq n$ \AND $m > 0$ \AND $s(m) = n$}
        \STATE $\mathrm{total} \gets \mathrm{total} + n$
    \ENDIF
\ENDFOR
\RETURN $\mathrm{total}$
\end{algorithmic}
\end{algorithm}

\section{Complexity Analysis}

\begin{proposition}
The algorithm runs in $O(N^{3/2})$ time and $O(1)$ auxiliary space.
\end{proposition}

\begin{proof}
The outer loop iterates $N - 2$ times, each calling $s(\cdot)$ at most twice at cost $O(\sqrt{N})$ per call (by Lemma~1). Total: $O(N \sqrt{N})$. No arrays beyond loop variables are needed.

A divisor-sum sieve computes all $s(n)$ for $n \leq N$ in $O(N \log N)$ time and $O(N)$ space, yielding an improved $O(N \log N)$ algorithm.
\end{proof}

\section{Answer}

\[
\boxed{31626}
\]

\end{document}
